\documentclass[a4paper]{article}
% Shared preamble for the English probability and statistics documents.
% Based on the original course preamble. Compile with pdfLaTeX.
\usepackage[top=1.5cm,bottom=2cm,left=2cm,right=2cm]{geometry}
\usepackage[T1]{fontenc}
\usepackage{lmodern}
\usepackage[utf8]{inputenc}
\usepackage[english]{babel}

\usepackage{amsmath}
\usepackage{amssymb}
\usepackage{amsthm}
\usepackage{mathrsfs}
\usepackage{mathtools}
\usepackage{etoolbox}
\usepackage{tikz}
\usetikzlibrary{shapes.multipart,arrows.meta}
\usepackage{xcolor}
\usepackage{enumerate}
\usepackage[shortlabels]{enumitem}
\usepackage{hyperref}
\usepackage{version}
\usepackage{bbold}
\usepackage{fontawesome5}
\usepackage{fancyhdr}
\lfoot{}
\usepackage{natbib}
\usepackage{graphicx}
\usepackage{pdfpages}
\usepackage{authblk}
    
\usepackage{blindtext}
\usepackage{multicol}
\usepackage{collcell}
\usepackage{wrapfig}
\usepackage{blkarray}
\usepackage{cancel}
\usepackage{appendix}
\usepackage{float}
\makeatletter
\input{tex/assets/slashbox.sty}
\makeatother
\hypersetup{unicode,colorlinks,citecolor=red,filecolor=black,
  linkcolor=red,urlcolor=cyan}
\setlength{\columnsep}{1cm}
\pagenumbering{gobble}

\def\numtext#1{%
  \ifcase#1\or one\or two\or three\or four\or five\or six\or seven\or eight\or nine\or ten\or
  eleven\or twelve\or thirteen\or fourteen\or fifteen\or sixteen\or seventeen\or eighteen\or nineteen\or twenty\or Lots\fi}

\newcommand{\Byte}[1]{%
  \def\nodecontents{}%
  \foreach[count=\l] \k in {#1} {\xappto{\nodecontents}{\noexpand\nodepart{\numtext{\l}}\k}}%
  \node[rectangle split,rectangle split parts=8,rectangle split horizontal,draw] at (2,2){\nodecontents};}

\newcommand{\D}{{\mathrm{d}}}
\newcommand{\cE}{{\mathcal{E}}}
\newcommand{\cG}{{\mathcal{G}}}
\newcommand{\cX}{{\mathcal{X}}}
\newcommand{\cY}{{\mathcal{Y}}}
\newcommand{\cB}{{\mathcal{B}}}
\newcommand{\cF}{{\mathcal{F}}}
\newcommand{\cO}{{\mathcal{O}}}
\newcommand{\cS}{{\mathcal{S}}}
\newcommand{\cT}{{\mathcal{T}}}
\newcommand{\given}{\, \big| \,}
\newcommand{\cL}{{\mathscr{L}}}
\newcommand{\cD}{{\mathcal{D}}}
\newcommand{\cA}{{\mathscr{A}}}
\newcommand{\cP}{{\mathscr{P}}}
\newcommand{\cR}{{\mathscr{R}}}
\newcommand{\A}{{\mathcal A}}
\newcommand{\one}{\mathbb{1}}
\newcommand{\mathbbm}[1]{\mathbb{#1}}

\providecommand{\C}{}
\renewcommand{\C}{\mathbb C}
\newcommand{\R}{\mathbb R}
\newcommand{\Q}{\mathbb Q}
\newcommand{\N}{\mathbb N}
\newcommand{\T}{\mathbb T}
\newcommand{\F}{\mathbb F}
\renewcommand{\P}{\mathbb P}
\newcommand{\Z}{\mathbb Z}
\newcommand{\E}{\mathbb {E}}
\newcommand{\Pbb}[1]{\P\left( #1\right)}
\newcommand{\Ebb}[1]{\E\left[ #1\right]}

\newcommand{\lb}{\left (}
\newcommand{\rb}{\right )}
\newcommand{\lbb}{\left [}
\newcommand{\rbb}{\right ]}
\newcommand{\labs}{\left |}
\newcommand{\rabs}{\right |}
\newcommand{\lbrb}[1]{\lb #1 \rb}
\newcommand{\lbbrbb}[1]{\lbb#1\rbb}
\newcommand{\labsrabs}[1]{\labs#1\rabs}
\newcommand{\lbbrb}[1]{\lbb#1\rb}
\newcommand{\lbrbb}[1]{\lb#1\rbb}
\newcommand{\langlerangle}[1]{\langle#1\rangle}
\newcommand{\lbcurly}{\{}
\newcommand{\rbcurly}{\}}
\newcommand{\curly}[1]{\lbcurly#1\rbcurly}
\newcommand{\abs}{\labsrabs}
\newcommand{\ind}[1]{1_{#1}}
\newcommand{\indep}{\perp \!\!\! \perp}
\renewcommand{\d}{\, \mathrm{d}}

\newtheorem{theorem}{Theorem}
\newtheorem{problemCnt}{Lemma}
\theoremstyle{definition}
\newtheorem{problem}[problemCnt]{Problem}
\newtheorem{problembook}[problemCnt]{\faBook \,  Problem}
\newtheorem{problempencil}[problemCnt]{\faPencil* \,  Problem}
\newtheorem{problemstar}[problemCnt]{\faStar\,  Problem}
\newtheorem{problemstarpencil}[problemCnt]{\faStar \, \faPencil* \,   Problem}
\newtheorem{probleminteresting}[problemCnt]{\faSmile[regular] \,   Problem}
\newtheorem{solution}{Solution}

\newcommand{\fwcell}[1]{\makebox[\arraystretch\normalbaselineskip]{$#1$}}
\newenvironment{enumeratei}{\begin{enumerate}[(i)]}{\end{enumerate}}
\newenvironment{enumeratea}{\begin{enumerate}[(a)]}{\end{enumerate}}
\newenvironment{enumeraten}{\begin{enumerate}[\upshape 1.]}{\end{enumerate}}

% Compatibility helpers needed by a few files in the combined archive.
\AtBeginDocument{\noindent{\footnotesize English translation prepared with ChatGPT (OpenAI, GPT-6), 26 October 2026, for translation of the Bulgarian text into English and LaTeX formatting.
The Bulgarian originals remain the reference; translation may contain errors.}\par\smallskip}
\newcommand{\MissingGraphic}[1]{%
  \begin{center}\fbox{\parbox{0.82\linewidth}{\centering Missing source image: \texttt{#1}}}\end{center}}



\setlength{\emergencystretch}{3em}
\allowdisplaybreaks
\begin{document}
\noindent Software Engineering \hfill Test 1, 16.12.2023: solutions
\par\smallskip\hrule\smallskip\noindent
\textit{English transcription of all twelve handwritten pages. The original calculations and their inconsistencies are preserved; this is a translation, not a corrected solution set.}

\section*{Problem 1}
\begin{enumerate}
\item We use Bayes' formula. Drawing a yellow, green, or blue ball means rolling one, two, or three dice, respectively. Thus
\begin{align*}
\P(\text{green}\mid\text{sum }3)
&=\frac{\P(\text{sum }3\mid\text{green})\P(\text{green})}
{\sum_i\P(\text{sum }3\mid\text{colour }i)\P(\text{colour }i)}\\
&=\frac{\frac2{36}\frac5{11}}
{\frac16\frac4{11}+\frac2{36}\frac5{11}+\frac1{6^3}\frac2{11}}
\approx29.13\%.
\end{align*}
\item The probability that the largest result is at least $3$ is
\[
\P(\max\geq3)=1-\P(\text{all results}<3)
=1-\left(\frac13\right)^5\approx99.59\%.
\]
For the product to be divisible by $10$, let $A=\{\text{product divisible by }2\}$
and $B=\{\text{product divisible by }5\}$. Then
\begin{align*}
\P(A\cap B)
&=1-\P(\overline A\cup\overline B)\\
&=1-\P(\overline A)-\P(\overline B)+\P(\overline A\cap\overline B)\\
&=1-\left(\frac12\right)^5-\left(\frac56\right)^5
+\left(\frac13\right)^5\approx57.10\%.
\end{align*}
\item The expected payoff is
\begin{align*}
&\left[\left(\frac13\right)^2+\left(\frac23\right)^2\right]
\left[\frac{\binom22}{\binom62}\,2
+\frac{2\cdot4}{\binom62}\,1\right]
-\frac49\left[\frac12\cdot3\right]\\
&\hspace{4cm}=-\frac8{27}<0.
\end{align*}
For a fair game, replace the payment by $x$ BGN:
\[
\frac{10}{27}-\frac49\frac x2=0
\quad\Longrightarrow\quad \boxed{x=\frac53}.
\]
\end{enumerate}

\section*{Problem 2}
\begin{enumerate}
\item By the definition of conditional probability,
\[
\P(X>m+n\mid X\geq n)
=\frac{\P(X>m+n)}{\P(X\geq n)}
=\frac{\P(X>m+n)}{\P(X>n-1)}.
\]
For $X\sim\operatorname{Ge}(p)$, with probabilities $pq^k$, $k\geq0$,
\begin{align*}
\P(X>n)&=\sum_{k=n+1}^{\infty}pq^k
=pq^{n+1}(1+q+\cdots)\\
&=\frac{pq^{n+1}}{1-q}=q^{n+1}.
\end{align*}
Consequently,
\[
\frac{q^{m+n+1}}{q^n}=q^{m+1}=\P(X>m).
\]
\item The handwritten argument starts with
\[
\P(Y>m+n\mid Y>n+1)=\P(Y>m),
\]
and then writes, for $n\geq1$,
\[
\P(Y>m+n)=\P(Y>n-1)\P(Y>m).
\]
Its recurrence is
\begin{align*}
\P(Y>m+1)
&=\P(Y>0)\P(Y>m)\\
&=\P(Y>0)\big[\P(Y>0)\P(Y>m-1)\big]\\
&=\cdots=\big(\P(Y>0)\big)^{m+2}.
\end{align*}
Let $\P(Y=0)=p$. Then $\P(Y>0)=1-p=:q$, and the source writes
\begin{align*}
\P(Y=k)&=\P(Y>k-1)-\P(Y>k)\\
&=q^{k+1}-q^k=q^k(1-q)=pq^k.
\end{align*}
It concludes that $Y\sim\operatorname{Ge}(\P(Y=0))$.
\item Using the i.i.d.\ assumption as in the handwritten solution,
\begin{align*}
\P(\min\{X_1,\ldots,X_n\}\leq t)
&=1-\P(X_{\min}>t)\\
&=1-\P(X_1>t)^n,\\
\P(\max\{X_1,\ldots,X_n\}\leq t)
&=\P(\text{all }X_i\leq t)\\
&=\P(X_1\leq t)^n.
\end{align*}
\item Distinguish the five dice. The time until the first six on each die has
the geometric distribution with parameter $1/6$. Here use the convention
$\P(X=k)=pq^{k-1}$ for $k\geq1$.
Let $X_i$ be the first-six time of die $i$. Then $K=\max_iX_i$ and
\[
\P(K\leq n)=\P(X_1\leq n)^5.
\]
Now
\begin{align*}
\P(X\leq n)&=\sum_{k=1}^{n}pq^{k-1}
=\sum_{k=0}^{n-1}pq^k
=\frac{p(1-q^n)}{1-q}=1-q^n,\\
\P(K\leq n)&=(1-q^n)^5,\\
\P(K=n)&=\P(K\leq n)-\P(K\leq n-1)\\
&=\boxed{(1-q^n)^5-(1-q^{n-1})^5},
\qquad q=\frac56.
\end{align*}
\end{enumerate}

\section*{Problem 3}
\begin{enumerate}
\item As in the birthday problem, the probability that all $n$ bit vectors
of length $m$ are distinct is
\[
\boxed{\frac{\prod_{k=0}^{n-1}(2^m-k)}{(2^m)^n}}.
\]
The expected total number of ones is
\[
\E\operatorname{Bin}(nm,\tfrac12)=\frac{mn}{2}.
\]
The expected number of vectors containing fewer than $m/4$ zeros is
$\E\operatorname{Bin}(n,p)=np$, where
\[
p=\sum_{k=0}^{\lceil m/4\rceil-1}
\P\big(\operatorname{Bin}(m,\tfrac12)=k\big).
\]
\item By linearity of expectation, the expected number of equal pairs is
\[
\E\sum_{i\ne j}\mathbf1_{\{X_i=X_j\}}
=\binom n2\P(X_1=X_2)=\frac{\binom n2}{2^m}.
\]
For $p\ne1/2$, the corresponding calculation is
\begin{align*}
&\binom n2\sum_{k=0}^m\binom mk
p^kp^k(1-p)^{m-k}(1-p)^{m-k}\\
&\quad=\binom n2\sum_{k=0}^m\binom mk
(p^2)^k\big((1-p)^2\big)^{m-k}\\
&\quad=\boxed{\binom n2\big[p^2+(1-p)^2\big]^m}.
\end{align*}
Alternatively, two coordinates agree with probability $(1-p)^2+p^2$.
Independence across coordinates gives
$\P(X_1=X_2)=[p^2+(1-p)^2]^m$.
\item The handwritten inclusion--exclusion argument for the expected number
of distinct vectors is
\begin{align*}
&\E\big[n-\#\text{equal pairs}+\#\text{equal triples}
-\#\text{equal quadruples}+\cdots\big]\\
&\quad=\E\left[\sum_{k=1}^{\infty}(-1)^{k+1}
\sum_{\substack{\text{distinct}\\i_1,\ldots,i_k}}
\mathbf1_{\{X_{i_1}=\cdots=X_{i_k}\}}\right]\\
&\quad=\sum_{k=1}^{\infty}(-1)^{k+1}\binom nk
\sum_{j=0}^{m}\binom mj(p^k)^j\big((1-p)^k\big)^{m-j}.
\end{align*}
The source's last line is
\[
\sum_{k=1}^{\infty}(-1)^{k+1}\big(p^k+(1-p)^k\big)^m.
\]
(The binomial coefficient in the preceding line is absent from that last
handwritten line.)
\item We seek the least $n$ for which the probability that all vectors are
distinct is less than $1/2$.

\textbf{Method I.}
\begin{align*}
\prod_{k=0}^{n-1}\left(1-\frac{k}{2^m}\right)&<\frac12\\
\Longleftrightarrow\quad
\log\prod_{k=0}^{n-1}\left(1-\frac{k}{2^m}\right)&<-\log2\\
\Longleftrightarrow\quad
\sum_{k=0}^{n-1}\log\left(1-\frac{k}{2^m}\right)
&\approx-\frac{n(n-1)}{2^{m+1}}<-\log2.
\end{align*}
Here $\log(1+x)\approx x$ for small $x$. Thus
\[
n(n-1)>2^{m+1}\log2,
\qquad
n_{\min}\approx\frac{1+\sqrt{1+2^{m+3}\log2}}2
\approx\sqrt{2^{m+1}\log2}.
\]

\textbf{Method II.}
\[
\frac{(2^m)!}{(2^m-n)!}\frac1{2^{mn}}<\frac12.
\]
Use Stirling's approximation $n!\approx\sqrt{2\pi n}(n/e)^n$:
\begin{align*}
\frac{(2^m)!}{(2^m-n)!}\frac1{2^{mn}}
&\approx
\frac{\sqrt{2\pi2^m}(2^m/e)^{2^m}}
{\sqrt{2\pi(2^m-n)}((2^m-n)/e)^{2^m-n}}\frac1{2^{mn}}\\
&=\sqrt{\frac1{1-n/2^m}}
\left(\frac{2^m}{2^m-n}\right)^{2^m-n}
2^{mn}\frac1{2^{mn}}\frac1{e^n}.
\end{align*}
For $n\ll2^m$, the source continues
\begin{align*}
&\approx\left(\frac1{1-n/2^m}\right)^{2^m-n}\frac1{e^n}\\
&=\frac{(1-n/2^m)^n}{(1-n/2^m)^{2^m}}\frac1{e^n}
\approx(1-n/2^m)^n<\frac12.
\end{align*}
Therefore
\[
\log(1-n/2^m)<-\frac{\log2}{n}
\quad\Longrightarrow\quad
\frac n{2^m}>\frac{\log2}{n}
\quad\Longrightarrow\quad
n>\sqrt{2^m\log2}.
\]
\textit{Where did a $\sqrt2$ go?}
In both cases, $n_{\min}=\Theta(2^{m/2})$ as $m\to\infty$.
\end{enumerate}

\section*{Problem 4}
The natural strategy is to choose \textit{higher} for $1$ and $2$ and
\textit{lower} for $4$ and $5$; for $3$ either choice is equivalent.
For simplicity, take \textit{lower} for $3$.

Since the probabilities are uniform, we can simply count the combinations
leading to each winning state. The original table is reproduced below.
The final column was left empty in the handwritten notes.
\[
\begin{array}{c|rrrrr}
\text{number }m\backslash\text{move }n&1&2&3&4&5\\ \hline
1&1&3&9&28&\\
2&1&4&12&37&\\
3&1&4&12&37&\\
4&1&3&9&28&\\
5&1&2&7&21&\\ \hline
\text{probability of reaching move }n&
\frac55&\frac{16}{5\cdot4}&\frac{49}{5\cdot4\cdot4}
&\frac{151}{5\cdot4\cdot4\cdot4}&
\end{array}
\]
The coloured arrows in the original show how the counts in the next
column are obtained by adding counts from the admissible preceding
states.

If we play until move $k$, the expected payoff is the probability of
reaching that move multiplied by the payoff at move $k$.
See the notebook for a generalization.

\medskip
If stopping at any time is also allowed (which is the realistic scenario),
the optimal expected payoff at each state is
\[
\max\left\{\text{current payoff},\
\sum_{i,k}\P(\text{move to state }(i,k))\,
\text{optimal payoff at }(i,k)\right\}.
\]
The partially filled table in the handwritten solution is
\[
\begin{array}{c|rrrr}
A&1&2&3&4\\ \hline
1&&&1000&1000\\
2&&&750&1000\\
3&&&&1000\\
4&&&&1000\\
5&&&&1000\\
\text{loss}&&&&
\end{array}
\]
For example,
\begin{align*}
A[1,3]
&=\max\left\{100,\frac14\,1000+\frac14\,1000+\cdots\right\}
=1000,\\
A[2,3]
&=\max\left\{100,\frac14\,0+\frac34\,1000\right\}=750.
\end{align*}
The first entry inside the maximum is the payoff if we stop; the second
is the expected payoff if we continue. In the second example, rolling a
$1$ means losing. Continue in the same way. See the notebook for the
solution of the general case.
\end{document}
